\section{Introduction}

\label{sec:intro}

The rise of AI agents---autonomous systems that combine LLM reasoning with tool use and memory~\cite{wang2024survey,xi2023rise}---has created demand for frameworks that simplify agent construction and orchestration. Current approaches fall into two categories:

\begin{enumerate}[leftmargin=1.4em]
    \item \textbf{Code-first frameworks} (LangChain~\cite{langchain2023}, CrewAI~\cite{crewai2024}, AutoGen~\cite{wu2023autogen}): Flexible but opaque. Agent behavior is encoded in Python/TypeScript code, making it difficult to visualize control flow, debug failures, or modify workflows without deep understanding of the codebase.
    \item \textbf{No-code builders} (Zapier AI, Make.com): Accessible but limited. Support simple linear chains but lack the expressiveness for conditional branching, parallel execution, loops, and multi-agent coordination required for complex workflows.
\end{enumerate}

Hanzo Flow occupies the middle ground: a visual programming environment with the expressiveness of code-first frameworks and the accessibility of no-code builders.

\subsection{Design Principles}

\begin{enumerate}[leftmargin=1.4em]
    \item \textbf{Graphs, not chains}: Workflows are directed acyclic graphs (DAGs) with typed edges, supporting arbitrary branching, merging, and parallel execution.
    \item \textbf{Static verification}: Type checking and constraint validation catch errors before execution.
    \item \textbf{Observable execution}: Every node execution is logged with full input/output traces, enabling replay and debugging.
    \item \textbf{Incremental deployment}: Workflows can be tested node-by-node, branch-by-branch, or end-to-end before deployment.
\end{enumerate}

\subsection{Contributions}

\begin{enumerate}[leftmargin=1.4em]
    \item A typed dataflow graph formalism for agent workflows with static verification (\S\ref{sec:formalism}).
    \item A visual debugger with execution replay and branching exploration (\S\ref{sec:debugger}).
    \item An adaptive execution engine with parallel branching and model fallback (\S\ref{sec:engine}).
    \item Evaluation on benchmarks and user study (\S\ref{sec:evaluation}).
    \item Production deployment analysis (\S\ref{sec:production}).
\end{enumerate}

